\section{Background}
\label{sec:background}

\subsection{Attention as a memory-bound kernel}

Scaled dot-product attention over a sequence of length $n$ with head dimension
$d$ computes
\begin{equation}
  \label{eq:attention}
  \mathrm{Attn}(Q, K, V) = \mathrm{softmax}\!\left(\frac{QK^{\top}}{\sqrt{d}}\right) V,
\end{equation}
where $Q, K, V \in \mathbb{R}^{n \times d}$. The arithmetic cost is
$\Theta(n^{2}d)$, but the quantity that governs runtime on current hardware is
the volume of key and value data read from device memory. Tiled formulations
such as FlashAttention~\cite{dao2022flashattention} restructure the computation
so that each key tile is read once per query block, which lowers the constant
factor without changing the asymptotic traffic.

Writing $B$ for the block size, the matrix decomposes into
$\lceil n/B \rceil^{2}$ blocks, and the traffic is proportional to the number of
blocks actually fetched. Our aim is to reduce that count.

\subsection{Why fixed patterns underperform}

Let $S_{ij}$ denote the softmax mass assigned by query block $i$ to key block
$j$. A fixed pattern selects a set $\mathcal{P}$ of block pairs independent of
the input, so its retained mass is $\sum_{(i,j) \in \mathcal{P}} S_{ij}$. In
measurements across 400 prompts, the block carrying the largest mass for a
given query block lay outside a 2048-token sliding window 31 percent of the
time, and outside any dilated stride pattern we tested 19 percent of the time.
A pattern chosen before the prompt arrives cannot capture that variability.

\subsection{Threat to accuracy}

Skipping a block replaces its contribution with zero before normalisation. The
resulting output error is bounded by the discarded mass times the spread of the
value vectors, which motivates the criterion in
Section~\ref{sec:design}: bound the discarded mass, not the number of blocks.
